\documentclass[10pt,a4paper]{amsart}
\usepackage[margin=19mm]{geometry}
\usepackage{amsmath,amssymb,mathtools}
\usepackage[scr=rsfs,cal=euler]{mathalfa}
\usepackage{xfrac}
\usepackage[unicode,hidelinks,bookmarks=false]{hyperref}
\newcommand{\ie}{i.e.\ }
\newcommand{\cf}{cf.\ }
\newcommand{\resp}{respectively\ }
\newcommand{\Subsection}{\subsection}
\providecommand{\nobreakdash}{\nobreak\hbox{-}}
\setlength{\emergencystretch}{2em}
\multlinegap0pt
\usepackage{xcolor}
\usepackage{amssymb}
\newtheorem{lemma}{Lemma}
\newtheorem{theorem}{Theorem}
\newcommand{\R}{\mathbb{R}}
\newcommand{\qe}{\mathbb {Q}^{n+1}_{\varepsilon}}
\title{On the inverse mean curvature flow by parallel hypersurfaces in space forms}
\author{Alancoc dos Santos Alencar, Keti Tenenblat}
\date{Source: arXiv:2403.14782v2 (CC BY 4.0)}
\begin{document}
\maketitle

\noindent\emph{Source and licence.} On the inverse mean curvature flow by parallel hypersurfaces in space forms. Version arXiv:2403.14782v2 (CC BY 4.0), distributed by arXiv under the Creative Commons Attribution 4.0 International licence, http://creativecommons.org/licenses/by/4.0/, as recorded in the arXiv licence metadata for this exact version and shown on the abstract page https://arxiv.org/abs/2403.14782v2. Full licence notice: \texttt{licenses/arxiv-2403.14782-CC-BY-4.0.txt}.\par\smallskip

\noindent\textit{Editorial scope.} Counted unit: the article's theorem res.17 (source0.tex lines 360--379). The article's complete original proof of it is delivered with it (source0.tex lines 657--713, one \texttt{proof} environment of 57 lines, which the article places away from the statement). No mathematics is written, repaired or re-derived: the statement and the proof are the article's own lines, in the article's own notation, and no retained line is substituted or re-flowed.  Retained uncounted as the source-local context the proof needs: lines 211--231; lines 233--252; lines 274--286; lines 509--515.  Nothing was omitted from any retained span, so the package carries no omission record.  No retained line contains a citation, so the wrapper carries no bibliography and no bibliography entry.  No retained line contains a figure environment, an image-inclusion command or drawing code, so no image is delivered and the package declares no asset dependency.  Editorial formatting only, in addition: an AMS \texttt{amsart} preamble that restates the article's own declarations and macros used by the retained lines, namely the environments \texttt{lemma}, \texttt{theorem} on separate counters, together with the standard packages those lines need.  The retained environments print in this document as Lemma 1, Theorem 1 in order of appearance, whereas the article prints the same items with the numbering of its own full document.  The complete native source of the article as distributed in the arXiv e-print for this exact version is bundled unchanged as \texttt{sources/156454/source0.tex}.
Considering the following notation
\begin{eqnarray}\label{eq.p3}
C_{\varepsilon}(\mu)= \left\{\begin{array}{ll}
\cos \mu,	& \mbox{if} ~~ \varepsilon = 1, \\ 
1,	& \mbox{if} ~~ \varepsilon = 0, \\
\cosh \mu,	& \mbox{if} ~~ \varepsilon = -1,
\end{array} 
\right. 
\mbox{ and ~~} 
S_{\varepsilon}(\mu)= \left\{\begin{array}{ll}
\sin \mu,	& \mbox{if} ~~ \varepsilon = 1, \\ 
\mu,	& \mbox{if} ~~ \varepsilon = 0, \\
\sinh \mu,	& \mbox{if} ~~ \varepsilon = -1, 
\end{array} 
\right.  
\end{eqnarray}
 the parallel hypersurface $F^\mu(M^{n})$ in $\qe$ are given by 
\begin{eqnarray}\label{eq.p4}
F^\mu (x)=	C_{\varepsilon}(\mu)F(x)+S_{\varepsilon}(\mu)N(F(x)),
\end{eqnarray}
where $\mu \in \R$ and $x \in M^{n}$. Now, we recall some basic facts about the parallel hypersurface $F^\mu(M^{n})$.

\begin{lemma}\label{res.4}
\noindent $(i)$ The unit normal vector field $ {N}^{\mu}$ of $F^\mu (M)$ is given by
\[
 {N}^{\mu}(x)=-\varepsilon S_\varepsilon(\mu)F(x)+C_\varepsilon(\mu)N(F(x)). \]

\noindent $(ii)$ If $E_{1},\ldots, E_{n} \subset T_{x}M$ is an orthonormal basis of eigenvectors of the second fundamental form $II_{x}$ of $F(M)$ in $x \in M$ and $k_{1},\ldots, k_{n}$ are the principal curvatures of $F(M)$ in $x \in M$ with respect to $N$, then the first fundamental form is given by 
\[
{g}^{\mu}_{x}(E_{i},E_{j})=\left[C_\varepsilon(\mu)-k_{i}(x) S_\varepsilon(\mu)\right]^{2} \delta_{ij}.  \nonumber
\]

\noindent $(iii)$ The second fundamental form ${II}^{\mu}_{x}$ of $F^\mu$ at $x \in M$ is given by
\[
{II}^{\mu}_{x}(E_{i},E_{j})=\left[C_\varepsilon(\mu)-k_{i}(x) S_\varepsilon(\mu)\right] \left[\varepsilon S_\varepsilon(\mu)+k_{i}(x) C_\varepsilon(\mu)\right] \delta_{ij}. 
\]

\noindent $(iv)$  The principal curvatures ${k}_{i}^{\mu}(x)$, $i=1$, $2$,..., $n$, of $F^\mu$ at $x \in M$, with respect to ${N}^{\mu}$, are given by
\[
{k}_{i}^{\mu}(x)=\frac{\varepsilon S_{\varepsilon}(\mu)+k_{i}(x)C_{\varepsilon}(\mu)}{C_{\varepsilon}(\mu)-k_{i}(x)S_{\varepsilon}(\mu)}.
 \]
\end{lemma}

\begin{theorem}\label{res.16}
Let $F:M^{n} \longrightarrow \qe$ be a connected  oriented hypersurface with inward   unit normal vector field $N$, i.e., such that the mean curvature  satisfies $H(x)>0$, $\forall x \in M$. Then, $F(M)$ is an initial data of a solution to the IMCF by parallel hypersurfaces if and only if $F(M)$ is an isoparametric hypersurface.
Moreover, for such a hypersurface,   
 the solution to the IMCF by parallel hypersurfaces with initial data $F(M)$ is given by
\begin{eqnarray}\label{eq.p.113}
 {F}^{t}(x)= C_{\varepsilon}(\mu(t))F(x)+S_{\varepsilon}(\mu(t))N(x),
\end{eqnarray}
$\forall (x,t) \in M \times I$, where $0 \in I \subset \R$, $C_{\varepsilon}$ and $S_{\varepsilon}$ are the functions given by $(\ref{eq.p3})$ and $\mu : I \longrightarrow \R$ is a smooth real function satisfying $\mu(0)=0$ and
\begin{eqnarray}\label{eq.p.114}
\prod_{i=1}^{n}\left( \, C_{\varepsilon}(\mu(t))-k_{i}S_{\varepsilon}(\mu(t))\, \right) = e^{t}, ~~ \forall t \in I, 
\end{eqnarray}
where $k_{i},\, i=1,...,n$ are the principal curvatures of $F(M)$, associated to $N$. 
\end{theorem}

\begin{lemma}\label{res.23}
Let $F:M^{n} \longrightarrow \mathbb{H}^{n+1}$ be an oriented isoparametric hypersurface with  unit normal vector field $N$ such that the mean curvature $H$ of $F(M)$ satisfies $H(x)>0$, $\forall x \in M$. Let ${F}^t:M^{n}  \longrightarrow \mathbb{H}^{n+1} \subset \mathbb{L}^{n+2}$ be the solution to the IMCF by parallel hypersurfaces with initial condition $F$ defined on the maximal interval $I$. If $(0,+\infty) \subset I$ and 
$$
\lim\limits_{t \longrightarrow +\infty}\cosh(\mu(t))=+\infty,
$$ 
then ${F}^{t}(M)$ converges to the boundary $\partial\mathbb{H}^{n+1}$ when $t \longrightarrow +\infty$.
\end{lemma}

\begin{theorem}\label{res.17}
	Let $F:\mathbb{S}^m \times \mathbb{H}^m \longrightarrow \mathbb{H}^{n+1} \subset \mathbb{L}^{n+2}$ be a connected  isometric immersion of a cylinder in the hyperbolic space, with 
	inward  
	unit normal vector field $N$ and two distinct principal curvatures, $k_{1}>1$ and $k_{2}$, with the same multiplicity $m$, such that $k_{1}k_{2}=1$. Then, the solution to the IMCF by parallel hypersurfaces, with initial condition $F$, is given by
	\begin{eqnarray}
	{F}^{t}(x)=\cosh(\mu(t))F(x)+\sinh(\mu(t))N(x), \nonumber
	\end{eqnarray}
	where
	\begin{eqnarray}\label{eq.p.104}
	\cosh (2\mu(t))=\frac{-e^{t/m}+a\sqrt{q(t)}} {a^2-1}, 
	\qquad 
	\sinh (2\mu(t))=\frac{-a e^{t/m}+\sqrt{q(t)}} {a^2-1}
	\end{eqnarray}
where
 $$q(t)= e^{2t/m}+ a^2-1, \quad \mbox{ and }\quad a=\frac{k_1+k_2}{2}= \frac{H}{n}$$ 
 $\forall t \in \R$, and $x\in \mathbb{S}^m \times \mathbb{H}^m$, i.e., ${F}^{t}$ is an eternal solution.
 Moreover, for each $t\in\R$,  ${F}^{t}(\mathbb{S}^{m}\times \mathbb{H}^{m})$ is a cylinder in the hyperbolic space with prinicpal curvatures satisfying $k_1^tk_2^t=1$, which converges to an $m$-dimensional submanifold of $\mathbb{H}^{n+1}$ when  
 when $t \rightarrow -\infty$ and it converges to  
  $\partial\mathbb{H}^{n+1}$ when $t \rightarrow +\infty$.
\end{theorem}

\begin{proof}[{\bf{Proof of Theorem \ref{res.17}}}]
From Theorem $\ref{res.16}$, we have that 
	\begin{eqnarray}
	{F}^{t}(x)=\cosh(\mu(t))F(x)+\sinh(\mu(t))N(x), \nonumber
	\end{eqnarray}
$\forall (x,t) \in  \mathbb{S}^{m}\times \mathbb{H}^{m} \times I$,	where
%	\begin{eqnarray}
$\prod_{i=1}^{n}(\,\cosh(\mu(t))-k_{i}\sinh(\mu(t))\,)=e^{t}$. 
%. \nonumber\end{eqnarray}
Since,  we are assuming that the principal curvatures $k_{1}$ and $k_{2}$ have the same multiplicity $m$,  
using the fact that 
 $k_{1}k_{2}=1$, we have that 
	\begin{eqnarray}\label{eq.p.119}
	\frac{k_{1}+k_{2}}{2}\sinh(2\mu(t))=\cosh(2\mu(t))-e^{t/m},
	\end{eqnarray}
$\forall t \in I$. 
We introduce the notation $a=(k_1+k_2)/2$, where $a>1$ by definition.   
The square on both sides of $(\ref{eq.p.119})$ implies that $\cosh(2\mu(t))$ must satify the algebraic equation 
	\begin{eqnarray}\label{eq.r.2}
	\left((a^2-1\right)\cosh^{2}(2\mu(t))+2e^{t/m}\cosh(2\mu(t))-(a^2+e^{2t/m})=0.
	\end{eqnarray}
Since  $\mu(0)=0$, it follows from \eqref{eq.r.2} and  \eqref{eq.p.119}, that $\forall t \in \R$,
	\begin{eqnarray}\label{eq.r.1}
	\cosh (2\mu(t))=\frac{-e^{t/m}+a\sqrt{q(t)}} {a^2-1}, 
	\qquad 
	\sinh (2	\mu(t))=\frac{-a e^{t/m}+\sqrt{q(t)}} {a^2-1},
	\end{eqnarray}
where 	
\begin{eqnarray*}
    q(t):=e^{2t/m}+a^2-1, 
\end{eqnarray*}
$\forall t \in \R$. The explicit expressions for  $\cosh(\mu(t))$ and $\sinh(\mu(t))$ can be easily obtained from \eqref{eq.r.1}.
Therefore, ${F}^t$ is defined for $t\in \R$, i.e., the solution to the IMCF  is eternal. Moreover,  $k_{1}>1$ implies that 
$	\lim\limits_{t \longrightarrow +\infty}\cosh(\mu(t))=+\infty$.

	Consider an orthonormal frame $\{E_{1},\ldots,E_{m},E_{m+1},...,E_{2m}=E_{n}\}$ of principal directions on $F(\mathbb{S}^{m} \times \mathbb{H}^{m}) \subset \mathbb{H}^{n+1} \subset \mathbb{L}^{n+2}$, in $x \in \mathbb{S}^{m}\times \mathbb{H}^{m}$, such that $E_{1}, ..., E_{m}$ are tangent to $\mathbb{S}^{m}$ and $E_{m+1}, ...,E_{n}$ are tangent to $\mathbb{H}^{m}$. From Lemma $\ref{res.4}$, 
 	the first fundamental form is given by 
 \begin{eqnarray*}
	I^{t}_{x}(E_{i},E_{j})=
\left\{
\begin{array}{lcl}	
	\left(k_{1} \sqrt{q(t)}+\frac{1-k_{1}^{2}}{2}\right)\delta_{ij},  &\quad  \mbox{ if } &  1\leq i,j\leq m, \\
\left(k_{2} \sqrt{q(t)}+\frac{1-k_{2}^{2}}{2}\right)\delta_{ij},  \quad  & \mbox{ if } &  m+1 \leq i,j\leq n.
\end{array}
\right.
	\end{eqnarray*}
Hence, from the expression of $q(t)$,  for $1 \leq i,j \leq m$, and  since  $k_1k_2=1$, we get
	\begin{eqnarray}\label{eq.a.8}
	\lim \limits_{t \longrightarrow -\infty}I^{t}_{x}(E_{i},E_{j})=0, 
	\end{eqnarray}
while for  $m+1 \leq i,j\leq n$,   
\[
\lim \limits_{t \longrightarrow -\infty}I^{t}_{x}(E_{i},E_{j})=(1-k_2^2)\delta_{ij}. 
\]	
 Therefore, ${F}^{t}(\mathbb{S}^{m}\times \mathbb{H}^{m})$ converges to an $m$-dimensional submanifold of $\mathbb{H}^{n+1}$ when $t \rightarrow -\infty$ and  it follows from Lemma \ref{res.23}, that  it converges to the boundary $\partial\mathbb{H}^{n+1}$ when $t \rightarrow +\infty$. 

\end{proof}

\end{document}
